\renewcommand{\Anno}{2023}

\renewcommand{\NumeroProblema}{3}

\renewcommand{\Testo}{
La Scuola Sant'Anna ha deciso di costruire una grande piramide nella piazza antistante la sede principale. Si tratterà di una classica piramide a base quadrata, con quattro facce triangolari tutte uguali tra di loro. Il volume $V$ della piramide è stato già stabilito, ma le altre dimensioni sono ancora oggetto di discussione tra tre diverse scuole di pensiero.
\begin{enumerate}
\renewcommand{\labelenumi}{(\alph{enumi})}

\item  Un primo gruppo vorrebbe che i quattro spigoli laterali della piramide (quelli che vanno dal vertice in alto ai quattro vertici del quadrato di base) formassero angoli di $45^{\circ}$ con il piano base.

In tale ipotesi, determinare (in funzione di $V$) la lunghezza del lato del quadrato di base della piramide.

\item  Un secondo gruppo vorrebbe porre l'attenzione sulla superficie laterale $S$ della piramide, cioè la somma delle aree dei quattro triangoli.

Determinare (in funzione di $V$) l'insieme dei valori possibili per $S$.

\item  Un terzo gruppo, infine, sta pensando ad una piramide trasparente che contiene la suo interno una sfera colorata.

Determinare (in funzione di $V$) il massimo volume possibile per questa sfera interna.

\end{enumerate}




%\clearpage

%
}

\renewcommand{\Risposta}{
% SSSUP-BOZZA-AI: Risposta; autore=ChatGPT; data=2026-11-08; stato=DA-VERIFICARE
% SSSUP-PROVENIENZA: AI: elaborazione nuova dal testo
Indichiamo con $a$ il lato della base.
\begin{enumerate}
\item[(a)] $a=(3\sqrt2\,V)^{1/3}$.
\item[(b)] $S\in[\sqrt3\,(18V^2)^{1/3},+\infty)$.
\item[(c)] Il massimo volume della sfera è $\pi V/8$, realizzato quando l'altezza della piramide è $\sqrt2$ volte il lato di base.
\end{enumerate}
}

\renewcommand{\Aiutino}{
% SSSUP-BOZZA-AI: Aiutino; autore=ChatGPT; data=2026-11-08; stato=DA-VERIFICARE
% SSSUP-PROVENIENZA: AI: elaborazione nuova dal testo
Scrivi $V=a^2h/3$. Per la superficie laterale usa l'altezza di una faccia, $\sqrt{h^2+a^2/4}$. Per la sfera massima puoi cercare il centro sull'asse della piramide e imporre l'uguaglianza delle distanze dalla base e dalle facce laterali; poi ottimizza rispetto al rapporto $2h/a$.
}

\renewcommand{\SoluzioneUfficiale}{}

\renewcommand{\SoluzioniAlternative}{
% SSSUP-BOZZA-AI: SoluzioniAlternative; autore=ChatGPT; data=2026-11-13; stato=DA-VERIFICARE
% SSSUP-PROVENIENZA: rielaborazione redazionale AI da appunti di soluzione 2023, archivio privato Concorsi/SSSUP_2023/Soluzioni/SSSUP_2023_Sol.pdf, pagina 8
\TitoloSoluzioneAlternativaA
\textbf{Punto (c): una dimostrazione algebrica mediante il discriminante.}
Indichiamo con $a$ il lato del quadrato di base, con $h$ l'altezza della piramide e con $r$ il raggio della più grande sfera che essa contiene. Come osservato nella soluzione redazionale, per simmetria possiamo scegliere il centro della sfera sull'asse della piramide; la sfera massima è tangente alla base e alle quattro facce. Il suo raggio è quindi il rapporto fra tre volte il volume e l'area totale della superficie:
\[
r=\frac{3V}{a^2+S},
\qquad
S=2a\sqrt{h^2+\frac{a^2}{4}}
=\sqrt{a^4+\frac{36V^2}{a^2}},
\]
poiché $a^2h=3V$.

Poniamo $t=a^2>0$. Dalla relazione precedente ricaviamo
\[
\frac{3V}{r}-t=\sqrt{t^2+\frac{36V^2}{t}}.
\]
Elevando al quadrato e semplificando, si ottiene
\[
2rt^2-3Vt+12Vr^2=0.
\]
Perché esista un valore reale di $t$, il discriminante di questo trinomio deve essere non negativo:
\[
\Delta=9V^2-96Vr^3\ge0.
\]
Pertanto
\[
r^3\le\frac{3V}{32},
\qquad
\operatorname{Vol}(\text{sfera})=\frac{4\pi}{3}r^3\le\frac{\pi V}{8}.
\]
L'uguaglianza si realizza davvero: scegliendo una piramide con $h=\sqrt2\,a$, si ha
\[
r=\frac{a}{2\sqrt2},\qquad
V=\frac{\sqrt2\,a^3}{3},\qquad
r^3=\frac{3V}{32}.
\]
La condizione sul discriminante fornisce quindi non soltanto un limite superiore, ma il massimo esatto, senza derivare alcuna funzione.
}

\renewcommand{\Osservazioni}{}

\renewcommand{\Commenti}{}

% Aggiunta redazionale: la soluzione preesistente non viene modificata.
\providecommand{\SoluzioneRedazionale}{}
\renewcommand{\SoluzioneRedazionale}{
% SSSUP-BOZZA-AI: SoluzioneRedazionale; autore=ChatGPT; data=2026-11-08; stato=DA-VERIFICARE
% SSSUP-PROVENIENZA: AI: elaborazione nuova dal testo
Consideriamo la piramide regolare descritta nel testo, con vertice sopra il centro del quadrato di base. Siano $a>0$ il lato e $h>0$ l'altezza. Vale $a^2h=3V$.

\paragraph{Punto (a).}
La proiezione di uno spigolo laterale sul piano di base misura $a/\sqrt2$. L'inclinazione di $45^\circ$ impone $h=a/\sqrt2$, quindi $V=a^3/(3\sqrt2)$ e $a=(3\sqrt2 V)^{1/3}$.

\paragraph{Punto (b).}
L'altezza di ogni faccia triangolare è $\sqrt{h^2+a^2/4}$. Pertanto
\[
S=2a\sqrt{h^2+a^2/4},\qquad
S^2=a^4+\frac{36V^2}{a^2}.
\]
La derivata rispetto ad $a$ si annulla solo quando $a^6=18V^2$, dove si ha il minimo. Il valore minimo è $S_{\min}=\sqrt3(18V^2)^{1/3}$. Poiché $S$ è continua e tende a infinito sia per $a\to0$ sia per $a\to\infty$, tutti i valori maggiori o uguali a tale minimo sono realizzati.

\paragraph{Punto (c).}
Per un raggio fissato, l'insieme dei centri delle sfere contenute nella piramide è convesso e invariante per le simmetrie del quadrato. Mediando un centro con i suoi ruotati di $90^\circ$ otteniamo un centro sull'asse, con la stessa ammissibilità. Possiamo quindi cercare la sfera massima sull'asse.

Se il centro ha quota $z$, la distanza dalla base è $z$ e quella da ciascun piano laterale è $a(h-z)/\sqrt{4h^2+a^2}$. Il minimo delle due è massimo quando coincidono. Il raggio massimo per quella piramide è dunque
\[
r=\frac{ah}{a+\sqrt{a^2+4h^2}}
 =\frac{h}{1+\sqrt{1+(2h/a)^2}}.
\]
Poniamo $t=2h/a$ e $u=\sqrt{1+t^2}>1$. Poiché $V=a^3t/6$, si ha
\[
\frac{r^3}{V}=\frac34\frac{t^2}{(1+\sqrt{1+t^2})^3}
=\frac34\frac{u-1}{(u+1)^2}.
\]
La derivata dell'ultima frazione ha il segno di $3-u$, perciò il massimo si ha per $u=3$, cioè $h/a=\sqrt2$. Si ottiene $r^3=3V/32$ e quindi
\[
\operatorname{Vol}(\text{sfera})=\frac{4\pi}{3}r^3\leq\frac{\pi V}{8},
\]
con uguaglianza nella piramide appena individuata.
}
